\section{Case three: an agent that modifies its own research harness}\label{sec:C33:06}
The research agent performs experiments in a sandbox and proposes changes to its own harness. The mutable objects include prompt templates, search strategy, context management, memory access, and tool orchestration. The evaluator contains protected tasks and graders. The candidate may inspect its search evidence but not final evaluation artifacts. This configuration resembles the bounded mechanisms studied by RRSI, Harness\textsuperscript{2}, and AIDE\textsuperscript{2}, with important differences in feedback and update targets~\cite{rrsi2026,corpus261014224,root_aide2026}.

Separate three artifact classes: the executor under test, the improver that proposes and selects changes, and the measurement system that scores them. Store immutable snapshots of all three at each accepted update. A proposal that changes a test must be classified as a measurement intervention, not quietly accepted as an executor improvement. A discovered evaluator bug can be real and worth repairing, but results before and after the repair need a consistent rescoring policy.

Suppose the current agent repeatedly overflows context and drops experiment configurations. A candidate harness offloads raw logs and retains a typed experiment record containing hypothesis, code revision, data split, random seed, metric definition, and result artifact. This may preserve scientific continuity with fewer active tokens. Its success should be tested on tasks where a later decision requires an earlier configuration, including changed metrics and failed runs. A summary that retains only the best score can create selection bias even if its memory use is efficient.

The candidate should be evaluated against a fixed-summary baseline, raw-log retrieval, and a no-change harness under matched total resources. Include the cost of creating and maintaining the structured records. If the candidate receives a stronger external critic, that critic's contribution must be separated. If the improver selects among many candidates, the final held-out evaluation should be independent of the selection stream. The relevant claim is transferable state management, not the ability to search longer on the same tasks.

The protected-evaluator design here should not be imputed to every cited method. Harness\textsuperscript{2} uses execution contrasts and model judgment for its task-level edits rather than ground-truth verifier access; its limitation section explicitly allows shared model bias and failed transfer~\cite{corpus261014224}. When repeated promotion reuses a benchmark, REUSE instead proves an admission guarantee under restricted decision-only feedback and an independent initial sample~\cite{ch02_reuse2026}. Revealing raw protected traces to the proposer violates that particular information restriction. Independent confirmation, adaptive reuse, and proxy-guided editing are distinct experimental contracts.

Now consider model post-training on the collected trajectories. If the retained record exposes information that the behavior policy did not see at a prior decision, training on a naively flattened history can change the conditioning problem. MemoryWalker identifies exact and approximate ways to address compaction-induced inconsistencies~\cite{j2026memorywalker}. The research system must record transform lineage and behavior inputs during collection; a later attempt to reconstruct them from a final summary may be impossible. The same memory design that helps execution can therefore create invalid training data if its information boundaries are ignored.

